% ============================================================================
% Core Feature: Master Data Administration
% ============================================================================
\section{Master Data Administration}

\begin{ugModule}{Overview}

  \textbf{Purpose:}
  Configure the reference data that drives document
  categorization, workflow statuses, and access policies.

  \textbf{Who Can Access:}
  \ugRole{Admin}
\end{ugModule}

\subsection{Document Types}

Define the categories of documents your organization manages:

\begin{itemize}[leftmargin=18pt]
  \item Examples: SOP, Surat Keputusan (SK), Proposal,
    Report, Work Instruction, Memo
  \item Full CRUD via \ugMenu{Settings > Master Data >
    Document Types}
  \item Each type has a \ugField{Name} and optional
    \ugField{Description}
\end{itemize}

\begin{ugSteps}
  \item Navigate to \ugMenu{Settings > Master Data >
        Document Types}.
  \item Click \ugButton{+ Add New}.
  \item Enter the \ugField{Type Name} and optional
        \ugField{Description}.
  \item Click \ugButton{Save}.
\end{ugSteps}

\subsection{Document Statuses}

The six default statuses are pre-configured:

\begin{center}
  \ugStatus{ugMutedText}{Draft}\quad
  \ugStatus{ugWarning}{Review}\quad
  \ugStatus{ugInfo}{Approved}\quad
  \ugStatus{ugSuccess}{Active}\quad
  \ugStatus{ugAccent}{Archived}\quad
  \ugStatus{ugDanger}{Rejected}
\end{center}

Additional custom statuses can be created if your workflow
requires them. Custom statuses can be assigned colors and
are available in the document status dropdown.

\subsection{Organizational Units}

Define the organizational structure for document grouping:

\begin{itemize}[leftmargin=18pt]
  \item Examples: IT, HR, Finance, Legal, Operations,
    Procurement, Quality Assurance
  \item Documents are assigned to a unit for filtering,
    reporting, and organizational dashboards
  \item Units can be created, edited, and deactivated
    (but not hard-deleted if documents reference them)
\end{itemize}

\subsection{Access Control Lists (ACL Groups)}

Define reusable access rule groups that can be assigned to
documents for batch access management:

\begin{itemize}[leftmargin=18pt]
  \item Create named groups (e.g., ``Finance Team'',
    ``Executive Board'')
  \item Assign users to groups
  \item When a document is linked to an ACL group, all
    members automatically receive access
  \item Simplifies access management for frequently shared
    document sets
\end{itemize}

\subsection{User \& Role Management}

\begin{tabularx}{\textwidth}{@{} l X @{}}
  \toprule
  \textbf{Function} & \textbf{Description} \\
  \midrule
  View Users
    & List all users synced from Keycloak with their
      roles and status \\
  Assign Roles
    & Change a user's role (Admin, Manager, Reviewer,
      Staff) \\
  Sync Users
    & Manually trigger re-sync of user data from
      Keycloak \\
  Deactivate User
    & Disable a user's access without removing their
      account \\
  \bottomrule
\end{tabularx}

\begin{ugNote}
  User accounts are managed in Keycloak and synchronized
  to Insight Doc. Changes to names, emails, and passwords
  should be made in Keycloak. Only role assignments are
  managed directly in Insight Doc.
\end{ugNote}

\newpage
